\documentclass[12pt]{article}
\usepackage{amsmath, amssymb}
\usepackage{geometry}
\geometry{a4paper, margin=1in}
\usepackage{enumitem}

\begin{document}

\title{\textbf{CSE400 - Lecture 22} \\ \Large Inequalities and Tail Bounds}
\author{}
\date{}

\maketitle

\section*{1. Tail of a Random Variable}

\textbf{Definition}

The tail of a random variable is:
\[
P(X > x)
\]

\textbf{Why do we care about tails?}
\begin{itemize}
\item Jobs delayed more than 24 hours
\item Packets dropped due to full buffer
\end{itemize}

\textbf{Key Idea}
\begin{itemize}
\item Mean $\rightarrow$ average behavior (easy)
\item Tail $\rightarrow$ rare extreme events (hard)
\end{itemize}

\textbf{Why are tails hard to compute?}
\begin{itemize}
\item Requires summing many small probabilities
\end{itemize}

\section*{2. Tail Examples}

\textbf{Example 1: Jobs Across Servers}

\textbf{Problem}

Suppose jobs are distributed among servers at random.

Find:
\[
P(\text{a particular server gets } \geq k \text{ jobs})
\]

\textbf{Model}

\[
X \sim \text{Binomial}(n,p)
\]
\[
P(X \geq k) = \sum_{i=k}^{n} \binom{n}{i} p^i (1-p)^{n-i}
\]

\textbf{Observation}
\begin{itemize}
\item Jobs are assigned randomly
\item Most servers get few jobs
\item Occasionally, one server gets many jobs
\item This is a tail event (rare overload)
\end{itemize}

\textbf{Example 2: Poisson Arrivals}

\textbf{Problem}

Jobs arrive according to a Poisson process with rate $\lambda$ jobs/hour.

Find:
\[
P(X \geq k)
\]

\textbf{Model}

\[
X \sim \text{Poisson}(\lambda)
\]
\[
P(X \geq k) = \sum_{i=k}^{\infty} e^{-\lambda} \frac{\lambda^i}{i!}
\]

\textbf{Observation}
\begin{itemize}
\item Jobs arrive randomly over time
\item Most hours have moderate arrivals
\item Sometimes there is a burst of arrivals
\item This is a tail event (sudden spike in load)
\end{itemize}

\section*{3. Why Tail Bounds?}

\textbf{Exact Computation}
\begin{itemize}
\item Complicated, long sums
\end{itemize}

\textbf{Idea}
\begin{itemize}
\item Instead of exact value: Find an upper bound
\item Easier to compute
\item Still gives a guarantee
\end{itemize}

\textbf{Key Idea}
\begin{itemize}
\item Approximate safely instead of computing exactly
\end{itemize}

\section*{4. Running Example (Key Idea)}

Flip a fair coin $n$ times.

\textbf{Distribution}
\[
X = \text{number of heads}, \quad X \sim \text{Binomial}(n, \tfrac{1}{2})
\]

\textbf{Mean}
\[
E[X] = \frac{n}{2}
\]

\textbf{Question}
\[
P\left(X \geq \frac{3n}{4}\right)
\]

\section*{5. Markov’s Inequality}

\textbf{Key Idea}
\begin{itemize}
\item Uses only the mean
\item Large values are rare
\end{itemize}

\textbf{Statement}

For non-negative $X$:
\[
P(X \geq a) \leq \frac{E[X]}{a}
\]

\textbf{Proof (as presented)}

\[
E[X] = \sum_{x=0}^{\infty} x p_X(x)
\]

\[
\geq \sum_{x=a}^{\infty} x p_X(x)
\]

\[
\geq \sum_{x=a}^{\infty} a p_X(x)
\]

\[
= a \cdot P(X \geq a)
\]

\[
\Rightarrow P(X \geq a) \leq \frac{E[X]}{a}
\]

\textbf{Running Example}

\[
X \sim \text{Binomial}(n, \tfrac{1}{2}), \quad E[X] = \frac{n}{2}
\]

\[
P\left(X \geq \frac{3n}{4}\right) \leq \frac{n/2}{3n/4} = \frac{2}{3}
\]

\textbf{Observation}
\begin{itemize}
\item Does not involve $n$
\item Weak bound
\end{itemize}

\textbf{Real-Life Example}

\[
P(X \geq 50000) \leq \frac{10000}{50000} = 0.2
\]

\section*{6. Chebyshev’s Inequality}

\textbf{Key Idea}
\begin{itemize}
\item Uses mean and variance
\item Controls deviation from mean
\end{itemize}

\textbf{Statement}
\[
P(|X - \mu| \geq a) \leq \frac{\text{Var}(X)}{a^2}
\]

\textbf{Proof Idea}

Apply Markov to $(X-\mu)^2$:
\[
P((X-\mu)^2 \geq a^2) \leq \frac{E[(X-\mu)^2]}{a^2}
\]

\[
= \frac{\text{Var}(X)}{a^2}
\]

\textbf{Running Example}

\[
\mu = \frac{n}{2}, \quad \text{Var}(X) = \frac{n}{4}
\]

\[
P\left(X \geq \frac{3n}{4}\right)
\leq P\left(|X - \mu| \geq \frac{n}{4}\right)
\leq \frac{n/4}{(n/4)^2} = \frac{4}{n}
\]

\textbf{Key Insight}
\begin{itemize}
\item Bound decreases as $n$ increases
\item Better than Markov
\end{itemize}

\section*{7. Chernoff Bound}

\textbf{Core Transformation}

\[
P(X \geq a) = P(e^{tX} \geq e^{ta})
\]

Apply Markov:
\[
P(X \geq a) \leq \frac{E[e^{tX}]}{e^{ta}}
\]

\textbf{Final Form}

\[
P(X \geq a) \leq \min_{t>0} \left\{ \frac{E[e^{tX}]}{e^{ta}} \right\}
\]

\textbf{Key Insight}
\begin{itemize}
\item Holds for all $t$
\item Choose best $t$
\item Stronger bound (exponential decay)
\end{itemize}

\section*{8. Chernoff Bound (Lower Tail)}

\[
P(X \leq a) \leq \min_{t<0} \left\{ \frac{E[e^{tX}]}{e^{ta}} \right\}
\]

\section*{9. Chernoff Bound for Poisson}

\[
E[e^{tX}] = \sum_{i=0}^{\infty} e^{ti} e^{-\lambda} \frac{\lambda^i}{i!}
\]

\[
= e^{-\lambda} \sum_{i=0}^{\infty} \frac{(\lambda e^t)^i}{i!}
\]

\[
= e^{\lambda(e^t - 1)}
\]


\end{document}